\section{Theoretical Setup}\label{sec:theoretical-setup}

\subsection{Problem Setting}\label{sec:problem-setting}
Fix integers $N\geq1$, $q\geq0$, and $m\geq0$, a number $T>0$, and a nondegenerate compact interval
$\Theta\subseteq[-T,T]$. Let $U\supseteq\Theta$ be an open interval. We use the one-dimensional
Balcan--Nguyen--Sharma common-chain convention \citep{balcan2024structured}. Let
$\eta=(\eta_1,\ldots,\eta_q)$ be a common triangular chain on $U$ satisfying
\begin{equation}
\label{eq:setup-display-001}
\eta_j'(\theta)=P_j(\theta,\eta_1(\theta),\ldots,\eta_j(\theta)),\qquad 1\leq j\leq q,
\end{equation}
and let $Q_i(\theta,y_1,\ldots,y_q)$, $0\leq i\leq N$, be output polynomials. With total degree,
\begin{equation}
\label{eq:setup-display-002}
M:=\begin{cases}\max_{1\leq j\leq q}\deg P_j,&q\geq1,\\0,&q=0,\end{cases}
\qquad
\Delta:=\max_{0\leq i\leq N}\deg Q_i,
\end{equation}
define
\begin{equation}
\label{eq:setup-display-003}
F_i(\theta):=Q_i(\theta,\eta_1(\theta),\ldots,\eta_q(\theta)),
\qquad
\widetilde F:=(F_0,F_1,\ldots,F_N),
\qquad F:=(F_1,\ldots,F_N).
\end{equation}
The coordinate $F_0$ is deterministic and only the coefficients multiplying $F_1,\ldots,F_N$ are random.

Let $B:U\to\mathbb R^{(N+1)\times(N+1)}$ be a supplied polynomial matrix,
\begin{equation}
\label{eq:setup-display-004}
B_{rs}(\theta)=\sum_{\ell=0}^{m}b_{rs,\ell}\theta^\ell,
\qquad 0\leq r,s\leq N,
\end{equation}
and set $T_*:=\max\{1,T\}$. Its static coefficient-height certificate is
\begin{equation}
\label{eq:setup-display-005}
\widehat\Lambda_{B,T}:=\left(\sum_{r=0}^{N}\sum_{s=0}^{N}
\left(\sum_{\ell=0}^{m}|b_{rs,\ell}|T_*^\ell\right)^2\right)^{1/2}.
\end{equation}

For $R>0$ and $0<\kappa<\infty$, let $\mathcal D_{N,R,\kappa}$ be the class of Borel probability laws
on $\mathbb R^N$ having one Lebesgue density $f_\mu$ supported on $[-R,R]^N$ with
$0\leq f_\mu\leq\kappa$ almost everywhere. Coordinates may be arbitrarily correlated. Assume this class is
nonempty, write $\mathcal D:=\mathcal D_{N,R,\kappa}$, and set
\begin{equation}
\label{eq:setup-display-006}
A:=(2R)^N\kappa.
\end{equation}
The affine event and its section are
\begin{equation}
\label{eq:setup-display-007}
\phi_\alpha(\theta):=F_0(\theta)+\langle\alpha,F(\theta)\rangle,
\qquad
H_\theta:=\{a\in\mathbb R^N:F_0(\theta)+\langle a,F(\theta)\rangle=0\}.
\end{equation}
Whenever $F_j(\theta)\neq0$, write the nonpivot coordinates as
$\beta=(\beta_i)_{i\neq j}\in[-R,R]^{N-1}$ and define
\begin{equation}
\label{eq:setup-display-008}
T_j(\theta,\beta):=-\frac{F_0(\theta)}{F_j(\theta)}
-\sum_{i\neq j}\beta_i\frac{F_i(\theta)}{F_j(\theta)}.
\end{equation}
Let $\Psi_j(\theta,\beta)$ insert $T_j$ in coordinate $j$ and $\beta_i$ in every other coordinate.
The associated graph Jacobian is $|\det D\Psi_j|=|\partial_\theta T_j|$.

The anchor and closure are theorem-critical conditions. Under them, $F\neq0$ on $\Theta$, so define
\begin{equation}
\label{eq:setup-display-009}
\gamma_F(\theta):=\frac{F(\theta)}{\|F(\theta)\|_2},\qquad
\Gamma_{\mathrm{proj}}(F):=\sup_{\theta\in\Theta}\|\gamma_F'(\theta)\|_2.
\end{equation}
Finally, for positive-length intervals, define
\begin{equation}
\label{eq:setup-display-010}
C^{\mathrm{Pf}}_{\mathcal D}(F;\Theta):=
\sup_{\mu\in\mathcal D}\sup_{\substack{I\subseteq\Theta\ \mathrm{interval}\\|I|>0}}
\frac{\Pr_{\alpha\sim\mu}[\exists\theta\in I:\langle\alpha,F(\theta)\rangle=0]}{|I|},
\end{equation}
\begin{equation}
\label{eq:setup-display-011}
C^{\mathrm{aff}}_{\mathcal D}(F_0,F;\Theta):=
\sup_{\mu\in\mathcal D}\sup_{\substack{I\subseteq\Theta\ \mathrm{interval}\\|I|>0}}
\frac{\Pr_{\alpha\sim\mu}[\exists\theta\in I:\phi_\alpha(\theta)=0]}{|I|}.
\end{equation}

\subsection{Assumptions}\label{sec:assumptions}
\begin{assumption}[Static one-dimensional parameter regime]\label{assump:parameter-regime}
$N\geq1$, $q\geq0$, $m\geq0$, $T>0$, $R>0$, and $0<\kappa<\infty$ are finite fixed instance data;
$\Theta\subseteq[-T,T]$ is nondegenerate compact, $U\supseteq\Theta$ is open, and
$\mathcal D_{N,R,\kappa}$ is nonempty. All presentation degrees, matrix coefficients, and support parameters
are fixed before a law or interval is selected.
\end{assumption}

\begin{assumption}[Balcan common-chain presentation]\label{assump:balcan-common-chain}
The chain and output polynomials use the triangular equations and degree convention in the problem setting, with ambient
parameter dimension $p=1$. Thus $q$ is chain length, $M$ is chain degree, and $\Delta$ is output-function degree.
\end{assumption}

\begin{assumption}[Anchored static derivative closure]\label{assump:anchored-derivative-closure}
There is a fixed $j_*\in\{1,\ldots,N\}$ with $Q_{j_*}\equiv1$, and the supplied deterministic matrix obeys
\begin{equation}
\label{eq:setup-display-012}
\widetilde F'(\theta)=B(\theta)\widetilde F(\theta)\qquad(\theta\in U).
\end{equation}
Consequently $F_{j_*}\equiv1$ and $\|F(\theta)\|_2\geq1$. No projective-speed, transversality, or
swept-area bound is assumed here.
\end{assumption}

\begin{assumption}[Arbitrarily correlated capped joint laws]\label{assump:cube-density-laws}
The coefficient vector is drawn from an arbitrary $\mu\in\mathcal D_{N,R,\kappa}$. The only distributional
inputs are cube support and the one full-joint density cap; no independence, marginal-density, or
conditional-density condition is imposed.
\end{assumption}

\subsection{Formalized Goal}\label{sec:formalized-goal}
The theorem below proves, without an additional assumption, the exact anchored derivative-closure
coefficient-sweep goal. It includes the static certificate and homogeneous normalized-curve identity, the
equivalent measurable pivot-chart inequalities, the coordinate-free affine rate and capacity, the sharper
homogeneous rate and capacity, the exact affine-monic presentation/certificate and probability recovery, and
Counter-example 1 with selected-law probability $\epsilon/(4\delta)$ and positive-length ratio
$1/(4\delta)$, all-law upper coefficient $1/\delta$, and raw certificate and projective speed $1/\delta$,
with their distinct probabilistic and certificate origins made explicit. All probabilities are ordinary
probabilities for a fixed law, all interval
suprema precede the outer law supremum,
and all stated vector, operator, Frobenius, scalar, Lebesgue, and Hausdorff norms and measures are Euclidean.
The conclusion is confined to presentations satisfying the supplied anchor and derivative-closure certificate;
it makes no claim that an unrestricted raw Pfaffian presentation admits that certificate or any polynomial
presentation-format bound.
